% ----------------------------------------------------------
% Capítulo 3 — Desenvolvimento
% ----------------------------------------------------------
\chapter{Desenvolvimento}
\label{cap:desenvolvimento}

Este capítulo apresenta os requisitos da solução, o ambiente computacional e a sequência metodológica. A tradução da consulta em parâmetros foi desenvolvida em três versões sucessivas, que compartilham a API, a ferramenta de busca e o banco e são resumidas na Seção~\ref{sec:visao-versoes}. As versões foram avaliadas em três conjuntos de consultas, o de referência, com \res{principal}{geral}{ntotal} consultas, o de validação, com \res{lotedados}{geral}{aceitas}, e o de teste, com \res{lote2dados}{geral}{aceitas} (Seção~\ref{sec:validacao}). Os resultados de cada fase estão no Capítulo~\ref{cap:resultados}.

% ============================================================
\section{Requisitos}
\label{sec:requisitos}
% ============================================================

Os requisitos da solução têm três origens. A primeira são as demandas apresentadas pelo 1º CGEO para este trabalho. A segunda é a análise do protótipo desse centro (Seção~\ref{sec:prototipo}), que fornece o esquema de busca a preservar e as limitações a superar. A terceira são as convenções que organizam o catálogo, a saber, a ET-PCDG \cite{etpcdg2016}, a articulação sistemática das folhas do SCN \cite{scn2016}, a divisão político-administrativa do IBGE e o Perfil de Metadados Geoespaciais do Brasil (Perfil MGB) \cite{perfilmgb2009}. O manual de anotação (Seção~\ref{sec:manual}) é a forma operacional dessas convenções e não constitui origem de requisito.

O Quadro~\ref{quadro:requisitos} lista os requisitos funcionais (RF) e não funcionais (RNF), com a origem de cada um. A tradução da consulta em parâmetros está desdobrada por funcionalidade, do RF2 ao RF7, para que cada ferramenta descrita na Seção~\ref{sec:res-modelagem} se ligue a um requisito determinado.

\begin{fixo}{quadro}
\centering
\caption{Requisitos da solução e a sua origem}
\label{quadro:requisitos}
\footnotesize
\begin{tabular}{|>{\centering\arraybackslash}p{1.2cm}|>{\raggedright\arraybackslash}p{9.0cm}|>{\raggedright\arraybackslash}p{4.4cm}|}
\hline
\textbf{Id.} & \textbf{Requisito} & \textbf{Origem} \\ \hline
RF1 & Receber a consulta em linguagem natural e devolver os parâmetros, os produtos e a resposta ao usuário & 1º CGEO; protótipo \\ \hline
RF2 & Reconhecer o lugar citado, seja a unidade da federação (UF), pelo nome ou pela sigla, seja o município, e restringir a busca pela interseção espacial com ele & protótipo; divisão político-administrativa do IBGE \\ \hline
RF3 & Reconhecer a folha pedida pelo código MI ou INOM, em qualquer grafia, ou pelo nome & 1º CGEO; articulação sistemática das folhas; limitação 4 do protótipo \\ \hline
RF4 & Reconhecer a escala em qualquer grafia e as expressões qualitativas de escala & 1º CGEO; escalas do catálogo; limitação 4 do protótipo \\ \hline
RF5 & Reconhecer o tipo de produto, o CGEO (área de suprimento) e o projeto pelos nomes e apelidos usuais & protótipo; ET-PCDG; limitação 4 do protótipo \\ \hline
RF6 & Resolver expressões de tempo em relação à data de execução, que o \textit{prompt} informa ao modelo como data atual, e distinguir publicação de criação & protótipo; metadados do catálogo \\ \hline
RF7 & Atender pedidos de ordenação e de quantidade, sem acrescentá-los quando não pedidos & protótipo \\ \hline
RF8 & Recusar, com explicação ao usuário, o pedido que não é do acervo, sem recusar pedido legítimo & 1º CGEO; limitação 2 do protótipo \\ \hline
RF9 & Tratar o pedido sem nenhum critério representável nos parâmetros & protótipo \\ \hline
RF10 & Informar a busca legítima sem resultado, com os critérios usados & 1º CGEO; protótipo \\ \hline
RF11 & Registrar, por consulta, os parâmetros, as chamadas, os tempos e os erros & limitações 5 e 6 do protótipo \\ \hline
RNF1 & Executar os modelos localmente & 1º CGEO \\ \hline
RNF2 & Latência compatível com o uso interativo, sem limiar fixado \textit{a priori} & 1º CGEO \\ \hline
RNF3 & Compatibilidade com o esquema do protótipo e com os metadados do BDGEx (ET-PCDG e Perfil MGB) & protótipo; ET-PCDG \\ \hline
RNF4 & Reprodutibilidade das avaliações & limitação 6 do protótipo \\ \hline
RNF5 & Manutenibilidade, com testes automatizados e integração contínua & limitação 1 do protótipo \\ \hline
RNF6 & Troca de modelo e de abordagem por configuração & limitação 3 do protótipo \\ \hline
\end{tabular}
\fonte{Elaborado pelos autores.}
\nota{As limitações numeradas são as da Seção~\ref{sec:prototipo-limitacoes}.}
\end{fixo}

O RNF2 não tem limiar, porque a latência depende do \textit{hardware} e da carga do ambiente de implantação, ambos fora do escopo do trabalho. Por esse motivo, a latência é reportada por ambiente de execução (Seção~\ref{sec:hardware}) e serve apenas ao desempate (Seção~\ref{sec:regra-decisao}). O atendimento de cada requisito é discutido na Seção~\ref{sec:atendimento}.

% ============================================================
\section{Ambiente Computacional}
\label{sec:ambiente}
% ============================================================

Esta seção apresenta os componentes da solução e o \textit{hardware}, justificado pelos requisitos dos modelos da Seção~\ref{sec:modelos}.

\subsection{Componentes da Solução e Tecnologias}
\label{sec:componentes}

A Figura~\ref{fig:componentes} mostra os componentes da solução e a tecnologia de cada um. A aplicação cliente envia a consulta à API, escrita em Python com o FastAPI \cite{fastapi2018} e o Pydantic \cite{pydantic2019}. O orquestrador da tradução, construído com o LangChain \cite{langchain2022}, liga as ferramentas ao modelo, servido pelo Ollama \cite{ollama2023} no formato GGUF \cite{gguf2023}, e recolhe as chamadas de ferramenta. As ferramentas auxiliares, usadas só na versão final, consultam a lista de municípios do IBGE, o índice de folhas do BDGEx e uma tabela de expressões de tempo. A ferramenta de busca confere o formato dos parâmetros e monta a consulta SQL parametrizada sobre o PostgreSQL com PostGIS \cite{postgis2001}, executado em contêiner, ambiente isolado e reprodutível de execução, gerido pelo Docker Compose. Na versão final, o validador confere ainda os valores antes da busca e devolve ao modelo os erros encontrados. O \textit{pipeline} de avaliação, fora do caminho do usuário, é usado para executar os conjuntos de consultas e pontuar as respostas.

\begin{fixo}{figure}
\centering
\begin{tikzpicture}[
  font=\scriptsize,
  comp/.style={rectangle, draw, rounded corners, align=center, minimum height=1.05cm, text width=3.7cm, inner sep=3pt, fill=gray!6},
  ext/.style={comp, fill=blue!7},
  dado/.style={comp, fill=green!8},
  aval/.style={comp, dashed, fill=white},
  arr/.style={-{Stealth[length=2mm]}, thick},
  darr/.style={{Stealth[length=2mm]}-{Stealth[length=2mm]}, thick}
]
  \node[ext] (cli) at (0,0) {\textbf{Aplicação cliente}\\consome a API (HTTP e JSON)};
  \node[comp] (api) [below=0.8cm of cli] {\textbf{API}\\Python, FastAPI; verificação das requisições com Pydantic};
  \node[comp, fill=orange!10, text width=4.2cm] (orq) [below=0.8cm of api] {\textbf{Orquestrador da tradução}\\Python, LangChain; laço de chamadas ao modelo (só na versão final), escrito no próprio código, com limites fixos};
  \node[comp, fill=purple!8, text width=3.4cm] (oll) [left=1.9cm of orq] {\textbf{Servidor de modelos}\\Ollama; pesos em formato GGUF, quantizados em 4 bits};
  \node[comp, text width=3.4cm] (fer) [right=1.9cm of orq] {\textbf{Ferramentas auxiliares do domínio}\\(só na versão final) Python; lista de municípios do IBGE, índice de folhas do BDGEx, tabela de expressões de tempo};
  \node[comp] (bus) [below=1.0cm of orq] {\textbf{Ferramenta de busca e validador}\\Python; conferência dos parâmetros (validador só na versão final) e montagem da SQL parametrizada (psycopg)};
  \node[dado] (db) [below=0.8cm of bus] {\textbf{Banco de dados}\\PostgreSQL 16 com PostGIS 3.4, esquema do protótipo, em contêiner Docker Compose};
  \node[aval] (pip) at (oll |- bus) {\textbf{\textit{Pipeline} de avaliação}\\Python; executa os conjuntos de consultas, pontua contra o gabarito e gera as tabelas};

  \draw[darr] (cli) -- node[right] {consulta / parâmetros, produtos e resposta} (api);
  \draw[darr] (api) -- (orq);
  \draw[darr] (oll) -- node[above, align=center, font=\tiny] {mensagens\\e chamadas} (orq);
  \draw[darr] (orq) -- node[above, align=center, font=\tiny] {argumentos\\e resultados} (fer);
  \draw[darr] (orq) -- node[right, align=left, font=\tiny] {parâmetros /\\produtos ou mensagem} (bus);
  \draw[darr] (bus) -- node[right, font=\tiny] {SQL / registros} (db);
  \draw[arr, dashed] (pip.north) -- node[left, align=right] {consultas\\dos conjuntos} (pip.north |- orq.south west) -- (orq.south west);
\end{tikzpicture}
\caption{Componentes da solução e tecnologia de cada um}
\label{fig:componentes}
\fonte{Elaborado pelos autores.}
\end{fixo}

Do protótipo foram reaproveitados o esquema do banco, os doze parâmetros de busca e a lógica das consultas SQL, portada sem o filtro pelo retângulo do mapa e sem a paginação, próprios da interface gráfica, que ficou fora do escopo. Foram reescritos em Python a API, antes em Node.js, a tradução, que passou da saída estruturada com o Instructor \cite{instructor2023} ao \textit{tool calling}, e o acesso ao banco. Foram eliminados o dicionário de normalização e o \textit{fallback} por expressões regulares, e cada falha do modelo passou a ser registrada e classificada (RF11). Foram criados as ferramentas auxiliares, o validador, os conjuntos de consultas e o \textit{pipeline} de avaliação.

O modelo não escreve SQL. Ele apenas preenche os parâmetros tipados de uma ferramenta fixa, com os quais a ferramenta de busca monta a consulta parametrizada (Seção~\ref{sec:tc-se-sql}). Adotou-se essa formulação porque o texto do usuário não chega à SQL, o que afasta a injeção \cite{pedro2025p2sql}, porque a geração de SQL tem acurácia ainda limitada em bancos reais \cite{li2023bird} e fica exposta à ambiguidade das perguntas diante do esquema \cite{liu2025survey}, e porque os enumerados da ET-PCDG permitem conferir cada valor antes da busca (RNF3). Em troca, a solução atende apenas às buscas que os doze parâmetros representam. Quanto à orquestração, o laço de chamadas da versão final é curto, tem limites fixos e foi escrito no próprio código do orquestrador, razão pela qual não se adotou o LangGraph \cite{langgraph2024}, que modela fluxos de agentes como grafos de estados.

A resposta ao usuário tem três desfechos, e em nenhum deles aparece produto que não esteja no banco:
\begin{enumerate}
  \item no pedido fora do acervo, não há busca, e a resposta é a frase de recusa dirigida ao usuário;
  \item na busca com produtos, a lista vem sempre do banco, e o modelo apenas a resume;
  \item na busca legítima sem produto (RF10), a API devolve a lista vazia e uma mensagem fixa com os critérios usados, que sugere retirar ou ampliar algum deles, sem nova chamada ao modelo.
\end{enumerate}

O modelo e a abordagem são escolhidos por configuração (RNF6). A API é empacotada em contêiner, e a sua configuração padrão é a versão final com o Gemma 4 E4B (Seção~\ref{sec:recomendacao}), com as versões anteriores ainda disponíveis. Os testes automatizados cobrem o \textit{schema}, as consultas SQL, as medidas e a API sem depender do Ollama nem do banco (RNF5), e a integração contínua os executa a cada mudança. Uma alternativa mais leve, com SpatiaLite ou DuckDB e execução direta dos modelos pelo Hugging Face, sem orquestração de contêineres, fica como trabalho futuro (Seção~\ref{sec:trabalhos-futuros}).

Como o protótipo não traz os dados do acervo, a demonstração ponta a ponta usou primeiro um banco sintético (27 UFs, 14 municípios, 5 áreas de suprimento e 117 produtos), e o acervo real foi carregado depois (Seção~\ref{sec:extracao}). O código, os conjuntos de consultas, o manual de anotação completo e os registros das rodadas estão no repositório deste trabalho \cite{correa2026pfcbusca}.

\subsection{\textit{Hardware} e Requisitos dos Modelos}
\label{sec:hardware}

Os modelos locais da Tabela~\ref{tab:modelos} ocupam de 2,5~GB a 7,1~GB em 4 bits, e o Phi-4 14B, modelo do protótipo, 9,1~GB, além da memória do contexto. Quando o total excede a memória da GPU (VRAM), parte do modelo é executada na CPU, e a latência aumenta. Esses requisitos justificam os três ambientes da Tabela~\ref{tab:ambientes}.

A estação local (Tabela~\ref{tab:ambientes}) representa o cenário de implantação modesta do RNF1 e do RNF2. Com 4~GB de VRAM, nenhum dos três modelos executados nela, o Gemma 4 E2B, o Gemma 4 E4B e o Qwen 3 4B, foi carregado inteiro na GPU, e a parcela carregada em cada caso é apresentada com os resultados da instalação (Seção~\ref{sec:res-instalacao}).

Na máquina virtual do Google Colab com GPU NVIDIA T4, o Ollama foi executado na própria máquina, sem serviço de modelo em nuvem. Com cerca de 15~GB utilizáveis, a GPU T4 comporta inteiro cada modelo local, inclusive o Phi-4 14B, o que permite distinguir o efeito do modelo do efeito da divisão entre GPU e CPU. Por isso as rodadas de qualidade dos modelos locais nos três conjuntos foram feitas na GPU T4, e a latência foi medida em cada ambiente.

Os modelos maiores em nuvem, GPT-OSS 20B e 120B \cite{openai2025gptoss} e Qwen 3.8 27B \cite{qwen2026qwen38}, servidos pelo Groq, excedem a VRAM da estação local e, executados fora dela, não atendem ao RNF1. A sua latência inclui a rede.

\begin{fixo}{table}
\centering
\caption{Ambientes de execução e uso de cada um}
\label{tab:ambientes}
\footnotesize
\begin{tabular}{|>{\raggedright\arraybackslash}p{2.3cm}|>{\raggedright\arraybackslash}p{3.9cm}|>{\raggedright\arraybackslash}p{4.0cm}|>{\raggedright\arraybackslash}p{4.0cm}|}
\hline
\textbf{Ambiente} & \textbf{\textit{Hardware}} & \textbf{Uso no trabalho} & \textbf{Justificativa} \\ \hline
Estação local & Intel Core i7-12700H; 32~GB de memória principal; NVIDIA GeForce RTX 3050 Laptop com 4~GB de VRAM; Windows 11 & latência e pares do teste A/B em parte do conjunto de referência; desenvolvimento da versão final & cenário de implantação modesta (RNF1, RNF2) \\ \hline
GPU T4 (Google Colab) & NVIDIA T4 com 16~GB nominais; Ollama na própria máquina virtual & rodadas de qualidade dos modelos locais nos três conjuntos; latência na GPU T4 & comporta cada modelo local inteiro, sem divisão entre GPU e CPU \\ \hline
Groq (nuvem) & \textit{hardware} do provedor, não controlado & primeira versão com os modelos maiores no conjunto de referência, incluídos os pares do teste A/B & só para situar os modelos locais; não atende ao RNF1 \\ \hline
\end{tabular}
\fonte{Elaborada pelos autores.}
\end{fixo}

% ============================================================
\section{Sequência Metodológica}
\label{sec:sequencia}
% ============================================================

A sequência metodológica segue cinco fases, que são a extração dos dados, a modelagem das funcionalidades, a instalação dos modelos, a validação dos modelos e a avaliação de performance. A Figura~\ref{fig:bpmn} apresenta essa sequência como mapa de processos na notação BPMN (\textit{Business Process Model and Notation}) \cite{omg2014bpmn} e mostra a iteração como ocorreu. A primeira versão foi modelada, os modelos foram instalados e validados no conjunto de referência, e só então os metadados do BDGEx foram extraídos para montar os conjuntos de validação e de teste. A segunda versão e a versão final foram modeladas em rodadas de correção e, congelada a versão final, a performance foi avaliada. Por fim, foi feita a análise complementar no acervo real. As tarefas destacadas em azul tiveram, além do apoio ao desenvolvimento do código, apoio de agentes de inteligência artificial (IA) (Seção~\ref{sec:uso-ia}). Depois da visão geral das versões, as subseções seguem as fases, com a extração dos dados (tarefas 1 e 5), a modelagem das funcionalidades (tarefas 2 e 6), a instalação dos modelos (tarefa 3), a validação dos modelos (tarefas 4 e 5) e a avaliação de performance (tarefas 7 e 8).

\begin{fixo}{figure}
\centering
\begin{tikzpicture}[
  font=\scriptsize,
  tarefa/.style={rectangle, draw, rounded corners=3pt, align=center, text width=2.75cm, minimum height=1.45cm, inner sep=2pt, fill=gray!6},
  ia/.style={tarefa, fill=blue!8},
  inicio/.style={circle, draw, minimum size=0.55cm, inner sep=0pt},
  fim/.style={circle, draw, ultra thick, minimum size=0.55cm, inner sep=0pt},
  porta/.style={diamond, draw, aspect=1, minimum size=0.95cm, inner sep=0pt, fill=yellow!12},
  fluxo/.style={-{Stealth[length=2mm]}, thick},
  nota/.style={font=\tiny, align=left, text width=3.0cm}
]
  \node[inicio] (ini) at (0,0) {};
  \node[tarefa] (t1) at (2.15,0) {1. Extrair os dados do protótipo e do IBGE};
  \node[tarefa] (t2) at (5.35,0) {2. Modelar as funcionalidades da primeira versão};
  \node[tarefa] (t3) at (8.55,0) {3. Instalar os modelos};
  \node[ia] (t4) at (11.75,0) {4. Validar os modelos no conjunto de referência e escolher o modelo};
  \node[ia] (t5) at (11.75,-2.7) {5. Extrair os metadados do BDGEx e montar os conjuntos de validação e de teste};
  \node[ia, minimum height=1.75cm] (t6) at (8.55,-2.7) {6. Modelar a segunda versão e a versão final};
  \node[porta] (g1) at (6.05,-2.7) {};
  \node[font=\tiny, align=center] at (6.05,-3.55) {versão final\\congelada?};
  \node[ia] (t7) at (3.25,-2.7) {7. Avaliar a performance};
  \node[tarefa, dashed] (t8) at (3.25,-5.0) {8. Análise complementar no acervo real};
  \node[fim] (fim) at (6.05,-5.0) {};
  \draw[->, thin] ([xshift=-0.12cm, yshift=0.14cm]t6.south) arc[start angle=200, end angle=-120, radius=0.11cm];

  \draw[fluxo] (ini) -- (t1);
  \draw[fluxo] (t1) -- (t2);
  \draw[fluxo] (t2) -- (t3);
  \draw[fluxo] (t3) -- (t4);
  \draw[fluxo] (t4) -- (t5);
  \draw[fluxo] (t5) -- (t6);
  \draw[fluxo] (t6) -- (g1);
  \draw[fluxo] (g1) -- node[above] {sim} (t7);
  \draw[fluxo] (g1.north) -- ++(0,0.75) -| node[pos=0.25, above] {não, nova rodada de correções} ([xshift=-0.6cm]t6.north);
  \draw[fluxo] (t7) -- (t8);
  \draw[fluxo] (t8) -- (fim);

  \node[nota, anchor=north] (n6) at ([xshift=0.3cm, yshift=-0.35cm]t6.south) {cada rodada mede a versão no conjunto de validação, agrupa os erros e define se a correção será escrita na descrição ou implementada no código};
  \draw[densely dotted] (t6.south) -- (n6.north);
  \node[nota, text width=2.0cm] (n7) at (0.55,-4.4) {conjunto de teste uma vez; outros modelos; teste A/B; comparação com o protótipo};
  \draw[densely dotted] (t7.south west) -- (n7.north);
  \node[rectangle, draw, fill=blue!8, minimum width=0.4cm, minimum height=0.25cm] (lg) at (0.0,-7.7) {};
  \node[font=\tiny, anchor=west] at (lg.east) {tarefa com apoio de agentes de IA, além do apoio ao código (Tabela~\ref{tab:uso-ia})};
\end{tikzpicture}
\caption{Sequência metodológica em notação BPMN}
\label{fig:bpmn}
\fonte{Elaborado pelos autores, na notação de \citeonline{omg2014bpmn}.}
\end{fixo}

\subsection{Visão Geral das Três Versões}
\label{sec:visao-versoes}

As três versões diferem apenas na tradução da consulta em parâmetros, e a API, a ferramenta de busca e o banco são comuns às três. Foram desenvolvidas em sequência, a primeira a partir do protótipo e cada uma das seguintes a partir dos erros da anterior, e a medida principal de cada versão foi feita num conjunto de consultas que não serviu para desenvolvê-la.

A primeira versão oferece ao modelo uma única ferramenta de busca, com os doze parâmetros do protótipo, e faz uma chamada ao modelo por consulta, sem novas tentativas. Difere do método do protótipo por usar \textit{tool calling} e por dispensar o dicionário de normalização e os exemplos resolvidos no \textit{prompt}. Para recusar um pedido, o modelo não chama a ferramenta. Essa versão foi usada no conjunto de referência para escolher o modelo e para a comparação com a saída estruturada e com o método do protótipo.

A segunda versão mantém a ferramenta de busca e uma única chamada ao modelo por consulta, inclui nas descrições dos parâmetros as convenções do manual de anotação (Seção~\ref{sec:manual}) e acrescenta uma ferramenta de recusa. Foi desenhada a partir dos erros da primeira versão no conjunto de referência e testada no conjunto de validação.

A versão final revisa as descrições e acrescenta quatro ferramentas auxiliares do domínio, que identificam nomes, normalizam códigos e escalas e resolvem expressões de tempo. Acrescenta também um pedido de esclarecimento, a verificação dos parâmetros com retorno ao modelo e a recusa contestada, pela qual a primeira recusa de uma consulta que cita um critério do catálogo volta ao modelo para confirmação, num laço de chamadas com limites fixos. Foi desenvolvida nos conjuntos de referência e de validação, congelada e, em seguida, testada uma única vez no conjunto de teste, com o procedimento descrito na Seção~\ref{sec:metodo-controle}.

A construção e o papel de cada conjunto estão na Seção~\ref{sec:validacao}, e o detalhe de cada versão, com a motivação de cada ferramenta, na Seção~\ref{sec:res-modelagem}. O Quadro~\ref{quadro:configuracoes} resume o que cada configuração avaliada contém.

\begin{fixo}{quadro}
\centering
\caption{Configurações avaliadas e o que cada uma contém}
\label{quadro:configuracoes}
\footnotesize
\setlength{\tabcolsep}{2.5pt}
\begin{tabular}{|>{\raggedright\arraybackslash}p{2.5cm}|>{\raggedright\arraybackslash}p{2.2cm}|>{\raggedright\arraybackslash}p{2.3cm}|>{\raggedright\arraybackslash}p{2.8cm}|>{\raggedright\arraybackslash}p{2.3cm}|>{\raggedright\arraybackslash}p{2.5cm}|}
\hline
\textbf{Configuração} & \textbf{Abordagem} & \textbf{Descrições dos parâmetros} & \textbf{Recusa} & \textbf{Ferramentas auxiliares, esclarecimento e laço} & \textbf{Verificação dos parâmetros} \\ \hline
Primeira versão & \textit{tool calling} & iniciais & não chamar a busca & não & não \\ \hline
Segunda versão & \textit{tool calling} & com as convenções do manual & ferramenta de recusa & não & não \\ \hline
Só com as descrições revisadas & \textit{tool calling} & revisadas & ferramenta de recusa & não & não \\ \hline
Com as ferramentas auxiliares & \textit{tool calling} & revisadas & ferramenta de recusa & sim & não \\ \hline
Versão final & \textit{tool calling} & revisadas & ferramenta de recusa, com recusa contestada & sim & sim, no meio da resposta \\ \hline
Controle & saída estruturada & revisadas & campo de recusa, com recusa contestada como nova tentativa & não & sim, como novas tentativas (no máximo duas) \\ \hline
\end{tabular}
\fonte{Elaborado pelos autores.}
\end{fixo}

\subsection{Extração dos Dados}
\label{sec:extracao}

A extração reuniu dados de três fontes, cujas estatísticas estão na Seção~\ref{sec:res-extracao}. Do protótipo \cite{cgeo2025prototipo} vieram os \res{principal}{e4b}{nconsultasP} casos de teste e o esquema do banco, com quatro tabelas. Do IBGE vieram a lista de municípios e as malhas das UFs e dos municípios. Do BDGEx foram coletados, pelo serviço CSW \cite{ogccsw2007}, em 28/09/2026, apenas os metadados públicos, sem baixar nenhum produto. Desses metadados foram extraídos os nomes de folha, os códigos, as escalas, os tipos de produto, os CGEOs e os anos usados nos conjuntos de validação e de teste, além do índice de folhas que as ferramentas da versão final consultam.

Para a análise complementar, os metadados foram carregados num PostgreSQL com PostGIS, no esquema do protótipo, junto das malhas do IBGE. Como o CSW não traz a geometria da folha nem campos de tipo, de escala e de projeto, a geometria de cada registro é a caixa envolvente informada, o tipo e a escala foram inferidos do título e do formato, e as áreas de suprimento, sem polígono oficial, foram aproximadas pela união das caixas dos registros de cada CGEO.

\subsection{Modelagem das Funcionalidades}
\label{sec:modelagem}

A modelagem definiu as perguntas e as exceções que o sistema deve resolver, fixou num manual a resposta correta de cada consulta e organizou o desenvolvimento em ciclos de correção.

\subsubsection{Tipos de Consulta e Exceções}
\label{sec:tipos-consulta}

O sistema deve resolver os tipos de consulta listados a seguir, que podem se combinar numa mesma consulta:
\begin{itemize}
  \item simples, como ``cartas de São Paulo'';
  \item compostas, como ``ortoimagens de Manaus em 1:50.000'';
  \item por código MI ou INOM, como ``INOM SG-22-X-A'';
  \item com tempo relativo, como ``nos últimos 3 meses'';
  \item com ordenação e quantidade, como ``as 3 cartas mais antigas'';
  \item informais ou ambíguas, com siglas e grafias sem acento.
\end{itemize}

Os tipos de consulta correspondem aos requisitos RF2 a RF7. Há também cinco exceções. As consultas fora do domínio, como ``cartas de Marte'', pedem o que o acervo não tem e devem ser recusadas (RF8). As consultas sem critério, como ``cartas do Nordeste'', não têm nenhum critério representável nos parâmetros e admitem a busca sem filtros ou a não busca (RF9). A busca legítima sem resultado deve ser informada ao usuário (RF10). O valor que os parâmetros não representam, como ``escala 1:10.000.000'', não tem comportamento correto determinado pelo texto e é um caso observacional, fora da pontuação. O nome ambíguo, como São Paulo, estado ou capital, admite mais de uma leitura (RF2). As consultas que citam ao menos um critério representável nos parâmetros, isto é, que não são fora do domínio nem sem critério, são chamadas consultas do domínio. A expressão distingue o pedido que deve resultar em busca do pedido que deve ser recusado ou esclarecido, e não o assunto, pois a consulta sem critério também trata do acervo.

\subsubsection{Manual de Anotação}
\label{sec:manual}

O manual de anotação é o documento que define a resposta correta de cada consulta, chamada gabarito. Foi escrito pelos autores antes de se anotarem as consultas, com base no gabarito original dos casos de teste do protótipo e nas convenções do catálogo. Tem seis princípios (evidência explícita no texto, forma canônica, leituras múltiplas explícitas, tratamento das consultas fora do domínio, consistência e casos observacionais) e regras para cada um dos doze parâmetros. Entre elas estão as regras de tempo relativo, cada uma identificada por um nome e resolvida na data de execução.

Parte das consultas admite mais de uma resposta correta, como ``últimos 3 meses'' (noventa dias ou três meses de calendário) ou ``as mais recentes'' (por data de publicação ou de criação). O gabarito registra todas as leituras aceitas por meio de três mecanismos, a saber, o valor alternativo (mais de um valor aceito num parâmetro), o campo opcional (parâmetro que pode faltar) e a leitura estrutural alternativa (o mesmo trecho em parâmetros diferentes), cada um justificado por uma regra. Entre as leituras aceitas, uma é marcada como leitura preferencial. Fora dessas leituras, a comparação é estrita.

O manual foi estendido aos conjuntos de validação e de teste sem alterar nenhuma anotação do conjunto de referência, e o texto completo está no repositório deste trabalho \cite{correa2026pfcbusca}. O manual nunca foi passado ao modelo, pois a sua função é definir o gabarito. Parte das suas convenções foi escrita nas descrições da ferramenta na segunda versão, e parte foi implementada em código, nas ferramentas e no validador da versão final.

\subsubsection{Ciclo de Desenvolvimento}
\label{sec:ciclo}

A segunda versão e a versão final resultaram de rodadas de um mesmo ciclo. A versão é medida num conjunto de desenvolvimento (o de referência ou o de validação), os erros são agrupados por tipo, e define-se o lugar de cada correção. A correção é implementada no código quando a regra é determinística e verificável na própria consulta, como a forma canônica, a tabela de datas ou as listas de UFs, municípios e folhas, e é escrita nas descrições e no \textit{prompt} quando se trata de orientar o uso das ferramentas. Na versão final, confere-se ainda que o validador não acusa nenhuma resposta certa do gabarito. Os ciclos da versão final foram registrados sobre uma amostra estratificada de \res{v3dev}{orientacoes}{n} consultas do conjunto de validação, com o Gemma 4 E4B na estação local, e agentes de IA foram empregados para classificar os erros e propor causas (Seção~\ref{sec:uso-ia}). O conjunto de teste ficou fora do ciclo (Seção~\ref{sec:papeis}).

\subsection{Instalação dos Modelos}
\label{sec:instalacao}

Os modelos locais foram selecionados por três critérios. O primeiro é o suporte nativo a \textit{tool calling} no Ollama, conferido pela etiqueta do catálogo oficial \cite{ollama2023}. O segundo é a existência de versões quantizadas em 4 bits executáveis em estações modestas, com parte do modelo na CPU quando a VRAM é insuficiente. O terceiro é o lançamento entre 2024 e 2026. Em conjunto, os modelos deviam cobrir mais de uma organização, e os quatro escolhidos cobrem três, a Alibaba, o Google e a Mistral AI, com dois portes da família Gemma 4. São o Qwen 3 4B \cite{qwen2025qwen34binstruct2507}, o Gemma 4 E2B e o Gemma 4 E4B \cite{team2026gemma4}, além do Mistral Nemo 12B \cite{mistral2024nemo} (Tabela~\ref{tab:modelos}). O Phi-4 14B \cite{abdin2024phi4} foi usado apenas na reprodução do método do protótipo (Seção~\ref{sec:metodo-prototipo}).

Todos foram instalados pelo Ollama com as etiquetas oficiais e executados com temperatura zero (o modelo escolhe sempre o \textit{token} mais provável), limite de 1.024 \textit{tokens} de saída e contexto de 4.096 \textit{tokens} na primeira e na segunda versão e de 8.192 na versão final, porque o \textit{prompt} desta, com várias ferramentas, se aproxima de 4.096 \textit{tokens}. Cada consulta é enviada sem conversas anteriores, e todos os modelos recebem, em cada versão, o mesmo \textit{prompt} de sistema, com a data de execução.

O raciocínio interno foi desligado por parâmetro, para que as latências fossem comparáveis, e o ganho possível com ele não foi medido. Na versão padrão do Qwen 3 4B, o raciocínio não podia ser desligado e, ao esgotar o limite de saída, impedia a chamada de ferramenta. Por isso, adotou-se a versão oficial sem raciocínio, Qwen 3 4B Instruct 2507, do mesmo porte. No Groq, usou-se o menor esforço de raciocínio aceito por cada modelo \cite{groq2026reasoning}, isto é, nenhum raciocínio no Qwen 3.8 27B e o esforço baixo nos GPT-OSS 20B e 120B, porque a família GPT-OSS não permite desligá-lo. O resultado da instalação está na Seção~\ref{sec:res-instalacao}.

\subsection{Validação dos Modelos}
\label{sec:validacao}

Nenhum conjunto de consultas foi usado para treinar ou ajustar os pesos de um modelo. Os conjuntos de referência e de validação orientaram o desenvolvimento das descrições, das ferramentas e do validador, e cada consulta é enviada ao modelo sozinha, sem exemplos resolvidos. Três nomes próximos designam coisas distintas. A validação dos modelos, fase desta seção, usou o conjunto de referência. O conjunto de validação, apesar do nome, serviu ao teste da segunda versão e ao desenvolvimento da versão final (Seção~\ref{sec:papeis}). O validador é o componente da versão final que confere os parâmetros antes da busca (Seção~\ref{sec:validador}).

\subsubsection{Conjunto de Referência}
\label{sec:conjunto-referencia}

O conjunto de referência reúne \res{principal}{geral}{ntotal} consultas de três origens disjuntas, cujas quantidades se somam. As consultas do protótipo são os \res{principal}{e4b}{nconsultasP} casos de teste do protótipo, rastreados até a linha do arquivo de origem, com a anotação original como leitura preferencial. As 40 consultas redigidas pelos autores cobrem lacunas, como as consultas simples, ausentes no protótipo, e casos de fronteira. As \res{principal}{geral}{nG} consultas geradas por modelos de frase combinam frases do português com listas fechadas do domínio (UFs, municípios, escalas, tipos, CGEOs, projetos, códigos, expressões de tempo e ordenações), com o gabarito conhecido por construção e geração determinística. As \res{principal}{geral}{nPN} consultas do protótipo e dos autores formam o \emph{subconjunto manual}.

O gabarito passou por uma auditoria em três passos. No primeiro, foram conferidas por checagens automáticas a unicidade, a validade dos valores, a resolução das regras de tempo em datas diversas, a rastreabilidade e a consistência, com o apoio de um anotador por regras escrito pelos autores. Um agente de IA de uma família que não está entre as avaliadas anotou as consultas de forma independente e às cegas, em dez blocos, só com o manual, a definição da ferramenta, a data de execução e o bloco. Os autores adjudicaram cada divergência e registraram a decisão. Um dos autores conferiu ainda à mão uma amostra estratificada de 40 consultas.

\subsubsection{Conjuntos de Validação e de Teste}
\label{sec:conjuntos-validacao-teste}

No conjunto de referência, as conclusões podem depender da regularidade das consultas geradas, a recusa é medida em poucas consultas fora do domínio, e o resultado pode refletir mais as descrições e o \textit{prompt} da versão avaliada que a abordagem. Por isso foram construídos dois conjuntos novos, segundo o mesmo procedimento. Dos metadados do BDGEx e da lista do IBGE foram sorteadas especificações de consultas a redigir, chamadas alvos (\res{lotedados}{geral}{alvos} para o conjunto de validação e \res{lote2dados}{geral}{alvos} para o de teste), distribuídas em oito famílias que seguem os tipos de consulta e as exceções da Seção~\ref{sec:tipos-consulta}, a saber, as consultas simples, as compostas, as por código, as de tempo relativo, as de ordenação, as informais ou ambíguas, as sem critério e as fora do domínio. Cada alvo fixa, antes de o texto existir, o gabarito segundo o manual, a forma como cada valor deve aparecer (como ``25k'' ou ``3o cgeo''), o que a consulta não pode dizer, um registro de linguagem e uma persona. Agentes de IA em sessões independentes redigiram as consultas a partir do alvo, isto é, dos valores a citar e da forma de escrevê-los, sem ver o objeto de parâmetros do gabarito nem o manual, e as repetições foram eliminadas por checagens automáticas. Outros agentes anotaram as consultas de forma independente, embaralhadas e sem o gabarito, e só foi aceita a consulta compatível com o gabarito nos dois sentidos (a leitura do anotador está entre as aceitas, e a preferencial está entre as que o anotador aceita). Uma amostra de 20\% recebeu uma segunda anotação, fora desse filtro. O filtro retém só as consultas cuja anotação independente é compatível com o gabarito e poderia excluir as formulações mais difíceis, mas as exclusões vieram quase todas das checagens automáticas e da eliminação de repetições. Dos \res{lotedados}{geral}{alvos} alvos do conjunto de validação e dos \res{lote2dados}{geral}{alvos} do de teste, restavam \res{lotedados}{geral}{dedup} e \res{lote2dados}{geral}{dedup} consultas antes do filtro e \res{lotedados}{geral}{aceitas} e \res{lote2dados}{geral}{aceitas} depois dele.

O conjunto de teste tem cerca de metade do tamanho do conjunto de validação, foi sorteado com outra semente aleatória e redigido em outras sessões de agentes, e nenhuma consulta dele foi usada no desenvolvimento. Os relatórios dos anotadores citaram trechos de algumas dezenas delas, e nada da versão final foi alterado a partir deles, nem há exemplo das descrições ou dos \textit{prompts} que coincida com uma consulta inteira desse conjunto. Nos dois conjuntos, porém, os nomes e os códigos foram sorteados das mesmas listas que as ferramentas consultam, e o alcance das ferramentas fora dessas listas não foi medido.

\subsubsection{Papéis dos Conjuntos e Congelamento}
\label{sec:papeis}

O conjunto de referência serviu à escolha do modelo, às comparações da primeira versão e ao desenvolvimento, e o de validação, ao teste da segunda versão e ao desenvolvimento da versão final. Com o de teste, a versão final foi medida uma única vez, porque um conjunto de teste usado para ajustar o sistema deixa de ser de teste \cite{dwork2015reusable}. Um número é dito \emph{dentro da amostra} quando foi medido num conjunto usado no desenvolvimento da própria versão medida, como a versão final no conjunto de referência, e é reportado apenas como registro, sem valor de teste. Uma análise é dita \emph{exploratória} quando não estava no plano registrado antes da rodada de teste e não alterou a versão final.

Para congelar a versão final antes da rodada de teste, calculou-se o resumo criptográfico (\textit{hash} SHA-256) do texto dos \textit{prompts}, das definições das ferramentas e do código e dos dados das ferramentas. O resumo é uma sequência curta calculada a partir do conteúdo, e qualquer alteração no conteúdo o modifica. A rodada de teste só foi executada depois de se conferir que os resumos eram iguais aos da versão congelada, e o plano de análise foi registrado no repositório antes dela. Para a reprodutibilidade (RNF4), cada rodada registra o modelo, as versões do \textit{software}, a configuração, a data de execução e o \textit{hardware}, e a pontuação é recalculada contra o gabarito vigente sem nova execução.

\subsection{Avaliação de Performance}
\label{sec:avaliacao}

A performance foi avaliada isoladamente e em comparação com o protótipo, com os conceitos estatísticos da Seção~\ref{sec:avaliacao-conceitos}.

\subsubsection{Medidas}
\label{sec:medidas}

A medida principal é a acurácia por consulta, ou acerto por parâmetros. A consulta está certa quando os parâmetros coincidem com uma leitura aceita, sem parâmetro a mais nem a menos, ou, nas consultas em que o correto é não buscar, quando não há busca. Os enumerados exigem igualdade exata, os nomes livres são comparados após normalizar caixa, acentos e pontuação, e os períodos exigem os mesmos limites. A acurácia é acompanhada do intervalo de confiança (IC) de 95\% de Wilson \cite{wilson1927probable}. No conjunto de referência, as rodadas da primeira versão com \textit{tool calling} nos modelos locais na GPU T4 foram executadas três vezes na mesma sessão, e as demais rodadas, em todos os conjuntos, uma vez. Nessas rodadas repetidas, as tabelas trazem a média das três repetições, e os testes estatísticos usam o resultado majoritário das três, porque repetições da mesma sessão quase sempre coincidem e contá-las como independentes exageraria a significância. Nos pares em que só uma das configurações foi repetida, como os da Tabela~\ref{tab:primeira-versao}, o resultado majoritário de uma é comparado com a única execução da outra.

Reportam-se também o F1 dos campos, ponderado pelas ocorrências, a acurácia no domínio e a recusa como classificação, cujo positivo é a consulta fora do domínio, com precisão, \textit{recall}, F1 da recusa e falsa recusa (não buscar numa consulta do domínio). Nas consultas sem critério, qualquer não busca conta como acerto. Essas consultas ficam fora da classificação da recusa, na qual conta como recusa qualquer resposta final sem busca, inclusive o pedido de esclarecimento não seguido de busca. A latência é medida depois de uma consulta de aquecimento e reportada pela mediana e pelo percentil 95 por ambiente. Ela abrange o intervalo do envio da consulta à resposta final, inclui as chamadas ao modelo, as ferramentas e o validador e exclui o tempo de execução da SQL.

Para verificar se os erros contados se deviam ao modelo ou ao gabarito, sortearam-se \res{auditoria}{erros}{n} execuções erradas da primeira e da segunda versão no conjunto de validação (sorteio à parte da amostra estratificada dos ciclos de desenvolvimento, Seção~\ref{sec:ciclo}), 40 por configuração de \textit{tool calling} e de saída estruturada, estratificadas por família de consulta, e um agente de IA as julgou segundo o manual. Essa auditoria dos erros é distinta da auditoria do gabarito descrita na Seção~\ref{sec:conjunto-referencia}.

\subsubsection{Regra de Decisão}
\label{sec:regra-decisao}

As medidas não são combinadas num escore único, e a acurácia por consulta é o critério principal. Duas configurações são comparadas nas mesmas consultas pelo teste de McNemar exato e bicaudal \cite{mcnemar1947note}, com uma observação por consulta. Há diferença quando o \textit{p}-valor é menor que \res{ab}{geral}{alfa}; do contrário, há empate, que não demonstra equivalência, porque diferenças de alguns pontos percentuais (p.p.) podem passar sem detecção. O \textit{bootstrap} pareado \cite{efron1993bootstrap} estima o tamanho da diferença, e a sobreposição de intervalos de Wilson não é usada como teste. Entre configurações empatadas, a preferência segue, sem teste, o F1 dos campos, a recusa e o custo (latência, chamadas ao modelo e parcela do modelo na GPU), critérios que não revertem uma diferença significativa de acurácia.

O plano registrado antes da rodada de teste fixou a medida, os testes, a comparação principal, os passos da versão final e a separação entre a resposta final e a primeira decisão, que é a primeira busca ou recusa anterior a qualquer retorno do validador. O nível de significância e a ordem dos desempates foram escritos depois das rodadas, e a correção de Holm \cite{holm1979simple}, aplicada às nove comparações planejadas do conjunto de teste, não estava no plano. As nove comparações são as do Gemma 4 E4B. Elas compreendem a comparação principal, as quatro diferenças entre passos consecutivos e as comparações da versão final e do controle com a melhor configuração anterior no conjunto de teste, a saída estruturada da segunda versão, e com uma arquitetura em duas etapas simulada a partir das próprias rodadas, na qual a decisão de buscar ou recusar é a da segunda versão com \textit{tool calling} e os parâmetros são os da saída estruturada da primeira versão.

\subsubsection{Comparação com o Protótipo}
\label{sec:metodo-prototipo}

O método de tradução do protótipo foi reproduzido em Python com o Phi-4 14B, o dicionário de 159 pares, o \textit{prompt} original em inglês com dez exemplos resolvidos, o \textit{schema} com campo de justificativa, até três novas tentativas e o \textit{fallback} por expressões regulares. As diferenças deliberadas são a temperatura zero (0,6 no protótipo), a exclusão da ordenação padrão que o protótipo acrescenta a toda busca e a ausência de conferência dos estados e municípios no banco. A saída é traduzida para os parâmetros pela correspondência usada no gabarito das consultas do protótipo. A reprodução foi executada uma vez, só no conjunto de referência, e comparada com a primeira versão nos quatro modelos locais e com a saída estruturada da primeira versão (Seção~\ref{sec:metodo-ab}).

\subsubsection{Passos da Versão Final e Controle}
\label{sec:metodo-controle}

No conjunto de teste, com o Gemma 4 E4B, a primeira e a segunda versão são seguidas de três passos cumulativos (Quadro~\ref{quadro:configuracoes}), aplicados em ordem fixa e designados como ``só com as descrições revisadas'', ``com as ferramentas auxiliares'' e ``com a verificação dos parâmetros''. O primeiro usa só as descrições revisadas, com uma chamada ao modelo. O segundo acrescenta as ferramentas auxiliares, o pedido de esclarecimento (respondido por um usuário simulado de resposta fixa, que diz não ter outros detalhes), o pedido para responder chamando uma ferramenta, enviado uma vez ao modelo quando a sua resposta não contém chamada, e o laço de chamadas. O terceiro acrescenta a verificação dos parâmetros e a recusa contestada, formando a versão final. Cada diferença mede o efeito do passo, mantidos os anteriores, e depende da ordem em que os passos foram aplicados. Nos três passos, a chamada de ferramenta que o modelo escreve no texto da resposta, em vez de emiti-la pelo canal de ferramentas, é interpretada como chamada (Seção~\ref{sec:versao-final}).

A versão final é comparada com um controle, a saída estruturada com as mesmas descrições revisadas e a mesma verificação dos parâmetros. Como a saída estruturada emite um único objeto e não consulta ferramentas durante a resposta (Seção~\ref{sec:tc-se-sql}), o validador e a recusa contestada são aplicados a ele como novas tentativas, no máximo duas, cada uma depois de um objeto completo. O controle não recebe as ferramentas auxiliares, o usuário simulado, o pedido de resposta por ferramenta, a leitura das chamadas escritas no texto nem a instrução do \textit{prompt} segundo a qual um lugar que o modelo não reconhece não justifica a recusa (Seção~\ref{sec:identificar-nome}). Se não houver diferença significativa entre o controle e a versão final, não se detecta efeito do mecanismo e das peças que só a versão final recebe, o que não demonstra que o ganho se deva apenas às descrições e à verificação. Se o controle tiver acurácia menor, a diferença mede o mecanismo de \textit{tool calling} junto com essas peças, sem separá-los, porque o controle que os separaria, com a saída estruturada recebendo o retorno das ferramentas auxiliares, não foi executado. Como as ferramentas e o validador codificam o manual que define o gabarito, inclusive a tabela de datas, a primeira decisão é analisada à parte da resposta final.

Conforme um plano registrado antes dessas rodadas, a comparação principal, o ganho sobre a primeira versão e a comparação entre \textit{tool calling} e saída estruturada na primeira versão foram repetidos com o Gemma 4 E2B e o Qwen 3 4B, na mesma versão final congelada. Essas rodadas conferiram os mesmos resumos criptográficos da versão congelada, mas usaram uma revisão posterior do restante do código e o Ollama 0.35.0, contra 0.34.4 nas do Gemma 4 E4B, e por isso as latências dos dois grupos não são comparáveis. O plano não previa comparar os modelos entre si.

\subsubsection{Teste A/B}
\label{sec:metodo-ab}

Para medir o ganho do \textit{tool calling} sobre a saída estruturada (Seção~\ref{sec:questoes}), as rodadas das duas abordagens foram reunidas num teste A/B. No \textit{tool calling}, as ferramentas definidas em JSON Schema são ligadas ao modelo pelo LangChain, e o modelo escolhe, pelo canal de ferramentas do Ollama ou do Groq, se chama cada uma e com quais argumentos. Na saída estruturada, o JSON Schema dos mesmos parâmetros é incluído no \textit{prompt}, e a resposta em modo JSON passa pela mesma classe Pydantic, que confere o formato sem restringir a geração. Na primeira versão, a saída estruturada não oferece meio de recusa, e a partir da segunda um campo próprio permite a recusa. A geração restrita ao \textit{schema} foi descartada, porque, com ela, os modelos encerravam o intervalo de datas logo depois do primeiro limite, sem o segundo, artefato que penalizaria a saída estruturada e é compatível com a observação de \citeonline{tam2024speak} de que restrições de formato alteram o desempenho dos modelos.

Cada par tem o mesmo modelo, as mesmas consultas, a mesma versão e o mesmo ambiente. São \res{ab}{geral}{npares} pares. Em \res{ab}{geral}{nparessemdesenv} deles, a versão comparada não foi desenvolvida no conjunto em que foi medida, e nos demais o resultado é dentro da amostra. O ganho em acurácia segue a regra de decisão, e o ganho em tempo é avaliado pela latência e pelo número de chamadas ao modelo, comparados só entre rodadas com a mesma GPU, a mesma versão do Ollama e a mesma parcela do modelo na GPU, o que exclui o Groq.

\subsubsection{Avaliação pela Execução no Acervo Real}
\label{sec:metodo-execucao}

O princípio da avaliação pela execução (Seção~\ref{sec:avaliacao-conceitos}), que compara o resultado da consulta, e não o seu texto, foi aplicado numa análise complementar, sem novas rodadas de modelos. No acervo carregado (Seção~\ref{sec:extracao}), cada leitura aceita das consultas do conjunto de teste em que a resposta correta é buscar foi expandida em buscas concretas (cada valor alternativo; cada campo opcional com e sem ele), com as datas resolvidas em 24/09/2026, a data de execução informada ao modelo nas rodadas do conjunto de teste (Seção~\ref{sec:resolver-periodo}), e a busca emitida por cada configuração nas rodadas existentes foi executada com a mesma SQL. O acerto por resultado exige o mesmo conjunto completo de cartas de alguma busca aceita ou, numa consulta sem critério, a não busca. Como muitas buscas do gabarito devolvem zero cartas, o acerto por resultado é calculado também apenas nas consultas cujo gabarito tem cartas. A análise é limitada porque o CSW não traz o projeto, grava uma só data por registro e não tem polígono oficial da área de suprimento, e porque o tipo e a escala foram inferidos do título.

\subsection{Uso de Ferramentas de Inteligência Artificial}
\label{sec:uso-ia}

A família de modelos usada como apoio, o Claude, da Anthropic, não está entre os modelos avaliados. A Tabela~\ref{tab:uso-ia} discrimina o uso por fase. O código foi desenvolvido com apoio do Claude Code. Em nenhum processo a IA definiu o gabarito. No conjunto de referência, ele foi definido pelos autores segundo o manual, e a anotação por agente serviu para medir a concordância e apontar divergências; nos conjuntos de validação e de teste, foi fixado por construção, antes da redação. A auditoria dos erros (Seção~\ref{sec:medidas}) coube a um único agente, sem conferência humana, e uma só família de modelos redigiu, anotou e julgou (Seção~\ref{sec:validade}).

\begin{fixo}{table}
\centering
\caption{Uso de ferramentas de inteligência artificial por fase da sequência metodológica}
\label{tab:uso-ia}
\footnotesize
\begin{tabular}{|>{\raggedright\arraybackslash}p{2.7cm}|>{\raggedright\arraybackslash}p{3.0cm}|>{\raggedright\arraybackslash}p{4.2cm}|>{\raggedright\arraybackslash}p{4.3cm}|}
\hline
\textbf{Fase} & \textbf{Ferramenta} & \textbf{Finalidade} & \textbf{Controle do resultado} \\ \hline
Todas as fases (código) & Claude Code, com os modelos Claude Opus 4.6 e Claude Opus 5.5 & desenvolvimento do código da coleta, do \textit{backend}, das ferramentas, do validador, do \textit{pipeline} de avaliação e dos testes & testes automatizados, revisão pelos autores e repositório público \\ \hline
Extração dos dados (coleta, tarefas 1 e 5) e instalação dos modelos (tarefa 3) & nenhuma além do código & não se aplica & não se aplica \\ \hline
Modelagem das funcionalidades (tarefa 6; na tarefa 2, só o código) & agentes Claude & classificar os erros das rodadas de desenvolvimento e propor causas & números recalculados por código; causas não testadas registradas como hipóteses \\ \hline
Validação dos modelos (tarefa 4 e montagem dos conjuntos na tarefa 5) & agentes Claude & anotar às cegas o conjunto de referência; redigir e anotar às cegas os conjuntos de validação e de teste & concordância medida; divergências adjudicadas pelos autores; checagens automáticas; filtro de compatibilidade; segunda anotação de uma amostra; conferência humana de 40 consultas do conjunto de referência \\ \hline
Avaliação de performance (tarefas 7 e 8) & agentes Claude & julgar \res{auditoria}{erros}{n} execuções erradas do conjunto de validação na auditoria dos erros (um único agente, sem conferência humana); classificar os erros restantes; conferir as análises & números recalculados por código \\ \hline
\end{tabular}
\fonte{Elaborada pelos autores. Os agentes usaram os modelos Claude Opus 5.5 e Claude Fable 5.1 (Anthropic).}
\end{fixo}
